% Options for packages loaded elsewhere
\PassOptionsToPackage{unicode,linktoc=all}{hyperref}
\PassOptionsToPackage{hyphens}{url}
%
\documentclass[
]{article}
\usepackage{amsmath,amssymb}
\usepackage{lmodern}
\usepackage{iftex}
\ifPDFTeX
  \usepackage[T1]{fontenc}
  \usepackage[utf8]{inputenc}
  \usepackage{textcomp} % provide euro and other symbols
\else % if luatex or xetex
  \usepackage{unicode-math}
  \defaultfontfeatures{Scale=MatchLowercase}
  \defaultfontfeatures[\rmfamily]{Ligatures=TeX,Scale=1}
\fi
% Use upquote if available, for straight quotes in verbatim environments
\IfFileExists{upquote.sty}{\usepackage{upquote}}{}
\IfFileExists{microtype.sty}{% use microtype if available
  \usepackage[]{microtype}
  \UseMicrotypeSet[protrusion]{basicmath} % disable protrusion for tt fonts
}{}
\makeatletter
\@ifundefined{KOMAClassName}{% if non-KOMA class
  \IfFileExists{parskip.sty}{%
    \usepackage{parskip}
  }{% else
    \setlength{\parindent}{0pt}
    \setlength{\parskip}{6pt plus 2pt minus 1pt}}
}{% if KOMA class
  \KOMAoptions{parskip=half}}
\makeatother
\usepackage{xcolor}
\usepackage[top=3cm,left=3cm,right=3cm,bottom=3cm]{geometry}
\usepackage{graphicx}
\makeatletter
\def\maxwidth{\ifdim\Gin@nat@width>\linewidth\linewidth\else\Gin@nat@width\fi}
\def\maxheight{\ifdim\Gin@nat@height>\textheight\textheight\else\Gin@nat@height\fi}
\makeatother
% Scale images if necessary, so that they will not overflow the page
% margins by default, and it is still possible to overwrite the defaults
% using explicit options in \includegraphics[width, height, ...]{}
\setkeys{Gin}{width=\maxwidth,height=\maxheight,keepaspectratio}
% Set default figure placement to htbp
\makeatletter
\def\fps@figure{htbp}
\makeatother
\setlength{\emergencystretch}{3em} % prevent overfull lines
\providecommand{\tightlist}{%
  \setlength{\itemsep}{0pt}\setlength{\parskip}{0pt}}
\setcounter{secnumdepth}{5}
\ifLuaTeX
  \usepackage{selnolig}  % disable illegal ligatures
\fi
\usepackage{csquotes}
\IfFileExists{bookmark.sty}{\usepackage{bookmark}}{\usepackage{hyperref}}
\IfFileExists{xurl.sty}{\usepackage{xurl}}{} % add URL line breaks if available
\urlstyle{same} % disable monospaced font for URLs
\hypersetup{
  pdftitle={D3Q User Manual},
  hidelinks,
  pdfcreator={LaTeX via pandoc}}

\title{D3Q User Manual}
\usepackage{etoolbox}
\makeatletter
\providecommand{\subtitle}[1]{% add subtitle to \maketitle
  \apptocmd{\@title}{\par {\large #1 \par}}{}{}
}
\makeatother
\subtitle{User Manual and Quick Reference}
\author{}
\date{}

\makeatletter
\def\@xobeysp{\mbox{}\space}
\def\verbatim@font{\normalfont\ttfamily\raggedright}
%\def\enquote#1{{``#1''}}
\makeatother

%lp
\setmainfont{TeX Gyre Schola}
\setmathfont{TeX Gyre Schola Math}

\begin{document}
\begin{titlepage}
\begin{center}
\def\version{1.2.0}
\def\qeImage{images/qe.png}


  \includegraphics[width=5cm,height=2cm,keepaspectratio]{\qeImage}
  \hspace{2cm} 
  \includegraphics[width=5cm,height=2cm,keepaspectratio]{images/anharmonic.png}
  
	       \vspace{5.5cm}
  \Huge D3Q User Manual (v.\version) for QE 7.1 
  \vspace{5.5cm}

  \includegraphics[width=3cm,height=1.5cm,keepaspectratio]{images/sciences-sorbonne.png}
  \hspace{2cm} 
  \includegraphics[width=3cm,height=1.5cm,keepaspectratio]{images/impmc.png}
  \hspace{2cm} 
  \includegraphics[width=3cm,height=1.5cm,keepaspectratio]{images/cnrs.png}

\end{center}

\end{titlepage}

{
\setcounter{tocdepth}{3}
\tableofcontents
}
\hypertarget{foreword}{%
\section{Foreword}\label{foreword}}

The thermal2 suite of codes has been written starting in 2010 by Lorenzo
Paulatto{[}1{]}. It is based on a previous code developed by Michele
Lazzeri and Stefano de Gironcoli{[}2{]} which has been distributed with
the Quantum-ESPRESSO suite of codes.

\hypertarget{copyright}{%
\subsection{Copyright}\label{copyright}}

All the files are provided under the
\href{https://www.gnu.org/licenses/old-licenses/gpl-2.0.en.html}{GPL
license, v2} or newer and, when possible, under the
\href{https://cecill.info/licences/Licence_CeCILL_V2.1-fr.html}{CeCILL
license}.

\hypertarget{citing}{%
\subsection{Citing}\label{citing}}

We would greatly appreciate if when using the d3q code you cite the
following papers where the underlying theory is described in detail:

\begin{itemize}
\tightlist
\item
  All applications: L. Paulatto, F. Mauri, and M. Lazzeri, Phys. Rev.~B
  87, 214303 (2013)
\item
  Exact BTE solution: G. Fugallo, M. Lazzeri, L. Paulatto, and F. Mauri,
  Phys. Rev.~B 88, 045430 (2013)
\item
  Spectral functions: L. Paulatto, I. Errea, M. Calandra, and F. Mauri,
  Phys. Rev.~B 91, 054304 (2015)
\item
  Finite size effects: L. Paulatto, D. Fournier, M. Marangolo, M.
  Eddrief, P. Atkinson, M. Calandra. Phys. Rev.~B 101 (20), 205419
  (2020)
\end{itemize}

\hypertarget{table-of-contents}{%
\section{Table of contents}\label{table-of-contents}}

\begin{itemize}
\tightlist
\item
  \protect\hyperlink{d3q-manual}{D3Q manual}
\item
  \protect\hyperlink{foreword}{Foreword}

  \begin{itemize}
  \tightlist
  \item
    \protect\hyperlink{copyright}{Copyright}
  \item
    \protect\hyperlink{citing}{Citing}
  \end{itemize}
\item
  \protect\hyperlink{table-of-contents}{Table of contents}
\item
  \protect\hyperlink{capabilities-of-the-d3q-code}{Capabilities of the
  d3q code}
\item
  \protect\hyperlink{compiling-the-code}{Compiling the code}
\item
  \protect\hyperlink{preparing-a-d3q-calculation}{Preparing a D3Q
  calculation}

  \begin{itemize}
  \tightlist
  \item
    \protect\hyperlink{running-the-pwx-calculation}{Running the pw.x
    calculation}
  \item
    \protect\hyperlink{running-the-phx-calculation}{Running the ph.x
    calculation}
  \item
    \protect\hyperlink{example-pwx-and-phx-input-files}{Example pw.x and
    ph.x input files}

    \begin{itemize}
    \tightlist
    \item
      \protect\hyperlink{input-of-pwx}{Input of pw.x}
    \item
      \protect\hyperlink{input-of-phx}{Input of ph.x}
    \end{itemize}
  \end{itemize}
\item
  \protect\hyperlink{d3q-input-file}{D3Q Input File}

  \begin{itemize}
  \tightlist
  \item
    \protect\hyperlink{qpoints-or-grid-specifications}{QPOINTS or GRID
    specifications}

    \begin{itemize}
    \tightlist
    \item
      \protect\hyperlink{inputd3q-namelist}{\&inputd3q namelist}

      \begin{itemize}
      \tightlist
      \item
        \protect\hyperlink{mode-character-default-single}{mode
        (CHARACTER, default: \enquote{single})}
      \item
        \protect\hyperlink{prefix-character-no-default}{prefix
        (CHARACTER, no default)}
      \item
        \protect\hyperlink{outdir-character-default-espresso_tmpdir}{outdir
        (CHARACTER, default \$ESPRESSO\_TMPDIR)}
      \item
        \protect\hyperlink{d3dir-character-default-outdir-or-espresso_d3dir}{d3dir
        (CHARACTER, default outdir, or \$ESPRESSO\_D3DIR)}
      \item
        \protect\hyperlink{fildrho_dir-character-default-outdir-or-espresso_fildrho_dir}{fildrho\_dir
        (CHARACTER, default outdir, or \$ESPRESSO\_FILDRHO\_DIR)}
      \item
        \protect\hyperlink{fild3dyn-character-default-anh}{fild3dyn
        (CHARACTER, default: \enquote{anh})}
      \item
        \protect\hyperlink{ethr_ph-real-default-1d-8}{ethr\_ph (REAL,
        default: 1.d-8)}
      \item
        \protect\hyperlink{amass-real-array-default-same-as-pwx}{amass
        (REAL, array, default: same as pw.x)}
      \item
        \protect\hyperlink{fildrho-character-default-drho}{fildrho
        (CHARACTER, default: \enquote{drho})}
      \item
        \protect\hyperlink{first-last-only-integer-default-1-0}{first,
        last, only (INTEGER, default: 1, 0)}
      \item
        \protect\hyperlink{offset-step-integer-default-0-1}{offset, step
        (INTEGER, default: 0, 1)}
      \item
        \protect\hyperlink{restart-logical-default-true}{restart
        (LOGICAL, default: .true.)}
      \item
        \protect\hyperlink{max_seconds-no-default-in-seconds}{max\_seconds
        (no default, in seconds)}
      \item
        \protect\hyperlink{max_time-format-hhmmss}{max\_time (format
        hh.mmss)}
      \item
        \protect\hyperlink{nk1-nk2-nk3-3x-integer-default-same-as-pwx}{nk1,
        nk2, nk3 (3x INTEGER, default: same as pw.x)}
      \item
        \protect\hyperlink{k1-k2-k3-3x-integer-default-same-as-pwx}{k1,
        k2, k3 (3x INTEGER, default: same as pw.x)}
      \item
        \protect\hyperlink{degauss-real-default-same-as-pwx}{degauss
        (REAL, default: same as pw.x)}
      \item
        \protect\hyperlink{print_star-logical-default-true}{print\_star
        (LOGICAL, default: .true.)}
      \item
        \protect\hyperlink{print_perm-logical-default-false}{print\_perm
        (LOGICAL, default: .false.)}
      \item
        \protect\hyperlink{print_trev-logical-default-false}{print\_trev
        (LOGICAL, default: .false.}
      \item
        \protect\hyperlink{safe_io-logical-default-false}{safe\_io
        (LOGICAL, default: .false.)}
      \end{itemize}
    \item
      \protect\hyperlink{grid-and-qpoints}{GRID and QPOINTS}
    \item
      \protect\hyperlink{d3_debug-namelist}{\&d3\_debug namelist}
    \end{itemize}
  \end{itemize}
\item
  \protect\hyperlink{running-d3q}{Running D3Q}

  \begin{itemize}
  \tightlist
  \item
    \protect\hyperlink{note-on-pools-parallelisation}{Note on pools
    parallelisation}
  \end{itemize}
\item
  \protect\hyperlink{output-format}{Output format}
\item
  \protect\hyperlink{data-flow}{Data flow}
\item
  \protect\hyperlink{bibliography}{Bibliography}
\end{itemize}

\hypertarget{capabilities-of-the-d3q-code}{%
\section{Capabilities of the d3q
code}\label{capabilities-of-the-d3q-code}}

The d3q code computes the third derivative of the Density Functional
Theory ground-state energy with respect to three harmonic perturbations,
identified by their wavevectors q1, q2 and q3=-q1-q2. The code can use a
certain number of methods:

\begin{itemize}
\tightlist
\item
  Norm Conserving pseudopotentials
\item
  Local Density Approximation (LDA) and Generalized Gradient
  Approximation (GGA) functionals.
\item
  Insulators and Metals (i.e.~partial occupation of the electronic
  bands)
\end{itemize}

On the other hand, the codes \emph{does not} implements the following
features:

\begin{itemize}
\tightlist
\item
  Ultrasoft pseudopotentials and PAW datasets
\item
  Advanced functionals, e.g.~meta-GGA, Grimme vdW corrections and
  non-local vdW functionals
\item
  Hybrid functionals
\item
  LDA+U, or self-interaction correction
\item
  External electric field, via Berry phase or with a sawtooth potential
\item
  spin-polarized systems
\item
  non-collinear spin and spin-orbit interaction
\end{itemize}

Please be aware of these limitations! Adding any of these features would
take a huge amount of work for developing the theory, implementing the
code and validating the results. There is very little chance that we
will implement any of them (perhaps with the exception of spin-polarized
systems) in the near future. On the other hand, if you are willing to
provide a consequent amount of workforce, we are available to provide as
much assistance as possible.

In most cases, we have have been able to obtain meaningful results
despite these limitation, sometimes combining phonons computed with more
advanced methods with 3rd order calculation computed on the same systems
but with a simpler method, assuming that more advanced correction would
not affect the 3rd order too much. If you want to take this approach we
recommend that you validate the assumption on a simpler test, where the
entire calculation can be performed with the methods available in d3q.x.
Alternatively, you can use one of the compatible real-space codes.

\hypertarget{compiling-the-code}{%
\section{Compiling the code}\label{compiling-the-code}}

The d3q.x code is distributed in a package bundled with the thermal2
codes for computing phonon-phonon interactions. The package has to be
decompressed inside the main directory of a Quantum-ESPRESSO source
distribution (which is usually called espresso-X.Y.Z, where X, Y and Z
identify the version). Please be aware that each release of the D3Q code
is developed and tested on top of a specific version of Quantum-ESPRESSO
and it will probably not work on any other version: Take care to choose
the correct package.

In order to decompress the package move inside the espresso-X.Y.Z
directory and issue the command tar xzvf
/where/it/is/d3q-A.B.C-qeX.Y.Z.tgz then move to the D3Q directory and
type \enquote{make} to compile the code. D3Q will use the same
compilation parameters as Quantum-ESPRESSO, as specified in the file
make.inc generated by the configure script of QE. The make command will
build both the D3Q code and the thermal2 codes.

\hypertarget{preparing-a-d3q-calculation}{%
\section{Preparing a D3Q
calculation}\label{preparing-a-d3q-calculation}}

The d3q.x code works in combination with the ph.x code from the
Quantum-ESPRESSO distribution, which in turn has to be run on top of a
total energy calculation performed with the pw.x code. As it is
impractical to store the wavefunctions and the wavefunction
perturbations, the d3q code recomputes them from the ground-state charge
density (produced by pw.x) and its variation with respect to a harmonic
phonon perturbation (computed by ph.x). The latter is, per default, not
stored to disk, here we will see the correct workflow to prepare all the
quantities required for a d3q calculation.

\hypertarget{running-the-pw.x-calculation}{%
\subsection{Running the pw.x
calculation}\label{running-the-pw.x-calculation}}

There is no specific keyword for the pw.x calculation, just a couple of
notes:

\begin{itemize}
\tightlist
\item
  Be sure to use a very tight value for conv\_thr, we recommend 1.d-9 or
  smaller, down to 1.d-12. It is important to start from a ground-state
  solution which is as close as possible to the variational minimum in
  order to ensure that the assumptions of the perturbation theory hold
\item
  Check the convergence of the plane-wave kinetic energy cutoff
  (ecutwfc) against a phonon calculation, not just the total energy. For
  example, you can check the convergence of the optical modes at Γ =
  (0,0,0), which is rather inexpensive, down to 0.1 cm-1. Phonon
  calculation can converge much slower than the total energy, especially
  when using GGA functionals.
\end{itemize}

\hypertarget{running-the-ph.x-calculation}{%
\subsection{Running the ph.x
calculation}\label{running-the-ph.x-calculation}}

Per default, the ph.x code does not save the variation of the charge
density (δn) to a file. In order to save it you can use the
fildrho=\enquote{filename} keyword in the \&inputph namelist of phonon,
which will store δn to filename in the outdir directory. This can be
sufficient, but note that it cannot be used when doing a phonon
dispersion calculation, as the δn of different q-points would overwrite
each other. Furthermore, if you are doing a phonon dispersion
calculation on a regular grid, only the points in the irreducible wedge
of the Brillouin zone will be computed explicitly, while the others are
obtained using symmetry. However, when d3q.x computes a regular grid of
(q1, q2, q3) triplets only q1 can be in the irreducible wedge: q2 will
span all the grid and q3 may not even be in the grid at all! In order to
solve these two problems, you have to add the drho\_star keywords to the
\&inputph namelist , with the following syntax:

\begin{verbatim}
drho_star%open = .true.
drho_star%dir  = "directory_name"
drho_star%ext  = "file_extension"
\end{verbatim}

Which will use symmetry to rotate δn from the irreducible wedge to all
the q-points in the regular grid; δn will be saved to a different file
for each q-point, in the directory drho\_star\%dir. We refer you to the
phonon documentation for mode details about these variables.

All this mechanism is performed automatically by the code: you do not
need to worry about the details. However, if you are interested, the
file names will be prefix.drho\_star\%ext.qpoint1. Where prefix and
drho\_star\%ext are the values of the corresponding variables the input
of ph.x, and qpoint1 is an automatically generated extension which
depend on the fractional coordinates of the q-point. In addition to each
δn file, there will be a file with the \enquote{.pat} extension which
contains the phonon polarization vector (also known as phonon pattern or
phonon eigenvector). Furthermore, the code will generate an index file
called prefix.drho\_star\%ext.dfile\_dir which contains a list of the
q-points; this index will be used by d3q.x to find the δn it needs: if
you are moving these files, do not forget to copy it along with the
rest.

\hypertarget{example-pw.x-and-ph.x-input-files}{%
\subsection{Example pw.x and ph.x input
files}\label{example-pw.x-and-ph.x-input-files}}

These are two minimalist input files for pw.x and ph.x for a crystal of
Silicon, which will allow you to proceed with a d3q.x calculation.

\hypertarget{input-of-pw.x}{%
\subsubsection{Input of pw.x}\label{input-of-pw.x}}

\begin{verbatim}
 &control
    calculation='scf'
    prefix='silicon',
    pseudo_dir = '../',
    outdir='/tmp'
 /
 &system
    ibrav=  2,
    celldm(1) = 10.20
    nat=  2,
    ntyp= 1,
    ecutwfc = 24
 /
 &electrons
    conv_thr =  1.0d-12
 /
ATOMIC_SPECIES
 Si  0.0  Si.pz-rrkj.UPF
ATOMIC_POSITIONS
 Si 0.00 0.00 0.00
 Si 0.25 0.25 0.25
K_POINTS AUTOMATIC
 4 4 4  0 0 0
\end{verbatim}

\hypertarget{input-of-ph.x}{%
\subsubsection{Input of ph.x}\label{input-of-ph.x}}

\begin{verbatim}
phonon_disp_of_silicon
 &inputph
   outdir='/tmp',
   prefix='silicon',
   ldisp=.true.
   nq1=4,   nq2=4,   nq3=4,
   tr2_ph=1.0d-16,
    drho_star%open = .true.
    drho_star%ext  = 'drho',
    drho_star%dir  = './FILDRHO'
 /
\end{verbatim}

\hypertarget{d3q-input-file}{%
\section{D3Q Input File}\label{d3q-input-file}}

The d3q.x code reads its input from a file, and a few environment
variables. If a parameter is specified both in the input file and in a
variable, it is the value in the file that is used. The input file has
the following structure:

\begin{verbatim}
title_line
&inputd3q
  ...
/
\end{verbatim}

\hypertarget{qpoints-or-grid-specifications}{%
\subsection{QPOINTS or GRID
specifications}\label{qpoints-or-grid-specifications}}

{[}debug instructions{]} The title line can be any title you wish to use
for your calculation, it will be printed in output but it has no effect
on the results. After the title line, you have to specify the \&inputd3q
namelist, which contain all the parameters of the calculation. After the
namelist, and depending on the value of the parameters therein, you will
have to specify either a series of q-points or the dimension of the
regular grid that you want to compute.

\hypertarget{inputd3q-namelist}{%
\subsubsection{\&inputd3q namelist}\label{inputd3q-namelist}}

The namelist start with the \enquote{\&inputd3q} keyword and ends with a
\enquote{/} in an otherwise empty line. It can contain the following
variables:

\hypertarget{mode-character-default-single}{%
\paragraph{\texorpdfstring{mode (CHARACTER, default:
\enquote{single})}{mode (CHARACTER, default: ``single'')}}\label{mode-character-default-single}}

This variable specify the mode in which d3q will operate, it can take
the following values:

\begin{itemize}
\tightlist
\item
  \enquote{single}: compute the 3rd order dynamical matrix for a single
  triplet (q1, q2, q3) of q-points. You will have to specify q1 and q2
  points, on two separate lines, after the namelist (units 2π/alat); q3
  is just -(q1+q2).
\item
  \enquote{gamma-only}: compute the D3 matrix of (0, 0, 0)
\item
  \enquote{gamma-q}: compute the matrix at (0,-q,q), like the old d3.x
  code; you have to specify the point after the namelist (units of
  2π/alat)
\item
  \enquote{partial}: compute a full grid of triplets around a specific
  point q0: (q0, q0+q, q0-q) with q = k1 b1/nq1 + k2 b2/nq2 + k3 b3/nq3
  with ki = 0, 1, \ldots{} , nqi. After the namelist, you will have to
  specify the point (units of 2π/alat) and then, on a new line, the size
  of the grid nq1 nq2 nq3. This method is in principle useful, as to
  compute the scattering matrix of phonon q0 you only need this set of
  D3 matrices, which can be especially useful for limiting the
  computational cost . However, using these matrices requires a
  specialized code, which is not implemented: the thermal2 suite of
  codes expects a full grid. If you badly need to use this feature,
  please contact us.
\item
  \enquote{full}: compute a full grid of triplets (q1, q2, q3), where q1
  and q2 span the grid and q3 =-(q1+q2). Symmetry is aggressively used
  to reduce the total number of triplets to compute, which makes it
  difficult to estimate; it is usually smaller than the square of the
  number of points in the irreducible wedge of the BZ, i.e.~if the
  phonon calculation included N q-points, the D3 calculation will
  probably include less than N2 triplets.
\end{itemize}

If you want to use the D3 matrices to do linewidth or thermal transport
calculations with the thermal2 codes, you definitely have to use the
\enquote{full} mode, although the other can be useful for testing
convergence or in other special cases.

\hypertarget{prefix-character-no-default}{%
\paragraph{prefix (CHARACTER, no
default)}\label{prefix-character-no-default}}

Prefix must be the same as in the input of pw.x and ph.x

\hypertarget{outdir-character-default-espresso_tmpdir}{%
\paragraph{outdir (CHARACTER, default
\$ESPRESSO\_TMPDIR)}\label{outdir-character-default-espresso_tmpdir}}

Outdir must be the same as in the input of pw.x and ph.x. Its default
value is read from the ESPRESSO\_TMPDIR environment variable, the
current directory is used if the variable is not set.

\hypertarget{d3dir-character-default-outdir-or-espresso_d3dir}{%
\paragraph{d3dir (CHARACTER, default outdir, or
\$ESPRESSO\_D3DIR)}\label{d3dir-character-default-outdir-or-espresso_d3dir}}

If you wish to write the d3 temporary data to a different place than the
pw data, you can use d3dir. Its default value is read from the
ESPRESSO\_D3DIR environment variable, and will eventually fallback to
the value of outdir. A further sub-directory will be created which
depends on the coordinates of the q-points, it is hence safe to use the
same value of d3dir even when running several d3q.x calculations at the
same time and with the same prefix, as long as they work on different
triplets. Data in d3dir is normally deleted after each triplet is
computed. If for some reason you need this data, you will have to change
d3q.f90 to call d3\_reset with parameter cleanup = .false., and
recompile the code.

\hypertarget{fildrho_dir-character-default-outdir-or-espresso_fildrho_dir}{%
\paragraph{fildrho\_dir (CHARACTER, default outdir, or
\$ESPRESSO\_FILDRHO\_DIR)}\label{fildrho_dir-character-default-outdir-or-espresso_fildrho_dir}}

The directory where the files with the variation of the charge density
have been stored. For the code to work, this must be the same as the
value of drho\_star\%dir in the input of ph.x. Note that only files
inside d3dir will be opened read/write, files inside outdir and
fildrho\_dir will be opened read-only, which ensure that even in case of
a catastrophic crash you should not need to repeat the pw.x and ph.x
calculations.

\hypertarget{fild3dyn-character-default-anh}{%
\paragraph{\texorpdfstring{fild3dyn (CHARACTER, default:
\enquote{anh})}{fild3dyn (CHARACTER, default: ``anh'')}}\label{fild3dyn-character-default-anh}}

The prefix of the files containing the D3 matrices. Note that the d3q.x
code will create a lot of output files (one for each of the explicitly
compute triplets, and one for each triplet obtained by symmetry); we
recommend that you insert a directory name in this variable to keep your
working directory tidy. E.g. you can use fild3dyn=\enquote{OUTPUT/anh},
the \enquote{OUTPUT} directory will be created automatically.

\hypertarget{ethr_ph-real-default-1.d-8}{%
\paragraph{ethr\_ph (REAL, default:
1.d-8)}\label{ethr_ph-real-default-1.d-8}}

The threshold for solving the non-self consistent Sternheimer equation.
The default value should work for everyone. Please note that this
threshold is not the same as the variable tr2\_ph from phonon input. It
is related to the threshold of the conjugate-gradient solver of the
Sternheimer equation; it is printed in the output of phonon, i.e.~the
value of \enquote{thresh} marked in bold in the following snippet from
ph.x:

\begin{verbatim}
Representation #  5 mode #   5 

Self-consistent Calculation 
    ...
 iter #   7 total cpu time : 4.0 secs   av.it.:   7.4
 thresh= 1.076E-08 alpha_mix =  0.700 |ddv_scf|^2 =  4.004E-17
\end{verbatim}

You can check the last iteration of your phonon calculation and use a
similar value for your d3 input.

\hypertarget{amass-real-array-default-same-as-pw.x}{%
\paragraph{amass (REAL, array, default: same as
pw.x)}\label{amass-real-array-default-same-as-pw.x}}

With amass you change the mass of atomic species. You can change it
again later if you wish, no part of the DFPT calculation depends on
ionic mass. The default value is read from the pw.x restart file.

\hypertarget{fildrho-character-default-drho}{%
\paragraph{\texorpdfstring{fildrho (CHARACTER, default:
\enquote{drho})}{fildrho (CHARACTER, default: ``drho'')}}\label{fildrho-character-default-drho}}

fildrho must be the same as drho\_star\%ext in ph.x. When this variable
is used, the d3q.x code will automatically search the file it needs
inside the fildrho\_dir directory. When computing a grid of triplets
calculation this is the only possible way to supply the fildrho files.

Alternatively, when computing single triplets, you can use fild1rho,
fild2rho and fild3rho to specify the rho variation at q point 1, 2 and
3. In special cases, i.e.~q2=-q3 and q2=q3, you only need to specify
fild1rho and fild2rho. In some other case, i.e.~q1=-q2 and q1=q2, you
only need fild1rho and fild3rho. In the Gamma-only case only fild1rho
file is used.

\hypertarget{first-last-only-integer-default-1-0}{%
\paragraph{first, last, only (INTEGER, default: 1,
0)}\label{first-last-only-integer-default-1-0}}

If you only want to compute a subset of the q-point you can use these
two variables. The d3q.x code will then compute all the triplets between
first and last. If you set last to zero, all the triplet starting from
first will be computed. Setting \enquote{only=x} is a shortcut for
\enquote{first=x} and \enquote{last=x}.

\hypertarget{offset-step-integer-default-0-1}{%
\paragraph{offset, step (INTEGER, default: 0,
1)}\label{offset-step-integer-default-0-1}}

Another way to select a subset of the triplets is to use the offset and
step variables. In this case the code will compute one triplet every
step, starting from 1+offset. Note that offset must be between 0 and
step-1. You can use first and last together with offset and step, but
note that first and last will always refer to the global triplet number.

\hypertarget{restart-logical-default-.true.}{%
\paragraph{restart (LOGICAL, default:
.true.)}\label{restart-logical-default-.true.}}

When restart is true, the code will check existing fild3dyn files and
skip the calculation of any triplet that has already been computed. This
option can also be used to refine a grid calculation, e.g.~if you have
already computed a 2x2x2 grid and want to compute a 4x4x4 one, simply
restart with the finer grid and the same input parameters, and the
triplets from the coarser grid will be skipped automatically. Restart of
partially-computed q-points is currently implemented but disabled
because 1) it is still unreliable 2) it requires that all input/output
is safely written to disk, which makes the calculation much slower. You
can activate it modifying the code at your own risk!

\hypertarget{max_seconds-no-default-in-seconds}{%
\paragraph{max\_seconds (no default, in
seconds)}\label{max_seconds-no-default-in-seconds}}

\hypertarget{max_time-format-hh.mmss}{%
\paragraph{max\_time (format hh.mmss)}\label{max_time-format-hh.mmss}}

The maximum running time after which the code will stop, you can only
set one of the two. The code will only check this condition after
finishing to compute each triplet.

\hypertarget{nk1-nk2-nk3-3x-integer-default-same-as-pw.x}{%
\paragraph{nk1, nk2, nk3 (3x INTEGER, default: same as
pw.x)}\label{nk1-nk2-nk3-3x-integer-default-same-as-pw.x}}

\hypertarget{k1-k2-k3-3x-integer-default-same-as-pw.x}{%
\paragraph{k1, k2, k3 (3x INTEGER, default: same as
pw.x)}\label{k1-k2-k3-3x-integer-default-same-as-pw.x}}

\hypertarget{degauss-real-default-same-as-pw.x}{%
\paragraph{degauss (REAL, default: same as
pw.x)}\label{degauss-real-default-same-as-pw.x}}

You can change the k-points grid and smearing used in d3q, do it with
care and only if you understand what you are doing, or for testing
purpose. nkX is the grid size along direction X, kX can be 1 or 0, it
means shift or no shift. degauss is the smearing in Ry.

\hypertarget{print_star-logical-default-.true.}{%
\paragraph{print\_star (LOGICAL, default:
.true.)}\label{print_star-logical-default-.true.}}

\hypertarget{print_perm-logical-default-.false.}{%
\paragraph{print\_perm (LOGICAL, default:
.false.)}\label{print_perm-logical-default-.false.}}

\hypertarget{print_trev-logical-default-.false.}{%
\paragraph{print\_trev (LOGICAL, default:
.false.}\label{print_trev-logical-default-.false.}}

The code saves the D3 matrix of each triplet computed and of the star of
those triplet. It does not compute the permutations (q1,q2,q3)
-\textgreater{} (q2,q3,q1); etc\ldots{} (although some symmetry
operations may be equivalent to permutations) It also does not use
time-reversal symmetry to send q-\textgreater-q. You can control this
behavior with these variables. Please note that the subsequent codes in
the thermal2 require print\_star=.true. in order to work correctly.

\hypertarget{safe_io-logical-default-.false.}{%
\paragraph{safe\_io (LOGICAL, default:
.false.)}\label{safe_io-logical-default-.false.}}

If you set safe\_io to true, every file is closed and reopened after
each write. This can make the code much slower, but can solve some
problems with parallel filesystem not syncing or corrupted file when
trying to restart. As long as partial q-point restart is not implemented
and tested, this option should left to its default value.

\hypertarget{grid-and-qpoints}{%
\subsubsection{GRID and QPOINTS}\label{grid-and-qpoints}}

After the \&inputd3q namelist, the grid dimension or q-points
coordinates have to be specified. Here, are some examples taken from a
hexagonal system:

Using mode=\enquote{gamma-q}, compute the point M=(1/2, 0, 0):

\begin{verbatim}
 0.5 0.0 0.0
\end{verbatim}

Using mode=\enquote{single}, compute (M, X, -M-X)

\begin{verbatim}
 0.5 0.0 0.0
 0.333333 0.333333 0.0
\end{verbatim}

In the \enquote{partial} case, compute 4x4x1 a grid centered around Γ:

\begin{verbatim}
 0. 0.0 0.0
 4 4 1
\end{verbatim}

In the \enquote{full} case, compute the entire 4x4x1 grid

\begin{verbatim}
 4 4 1
\end{verbatim}

\hypertarget{d3_debug-namelist}{%
\subsubsection{\&d3\_debug namelist}\label{d3_debug-namelist}}

The variables inside the d3\_debug namelist can be used to disable some
parts of the calculation or to print more output. WARNING! no sanity
check is done on these variables, setting this namelist can result in
crash and random results. No bug report will be taken into consideration
if the code was run with debug variables! All variable in this namelist
default to .true.; except for dbg\_write\_d3\_part and dbg\_full\_bands
which default to false.

\begin{itemize}
\tightlist
\item
  dbg\_do\_dwfc: compute d\^{}q psi
\item
  dbg\_do\_dpdvp: pre-compute \textless d psi\textbar{} dV\textbar{}
  psi\textgreater{}
\item
  dbg\_do\_dpdvdp: compute the \textless d psi\textbar{} d V\textbar{} d
  psi\textgreater{} term
\item
  dbg\_do\_dpdpdv: compute the \textless d psi\textbar d
  psi\textgreater\textless psi\textbar{} dV \textbar psi\textgreater{}
  term
\item
  dbg\_do\_drhod2v: compute the d rho d\^{}2 V term (and related wfc
  terms)
\item
  dbg\_do\_rhod3v: compute the tho d\^{}3 V term (and related)
\item
  dbg\_do\_ion: compute Ion-Ion terms (using the Ewald sum technique)
\item
  dbg\_do\_smearing: compute metal-only terms
\item
  dbg\_do\_smr\_ijk: metal-only term not involving Fermi-energy shift
\item
  dbg\_do\_smr\_ij: metal-only term involving one Fermi-energy shift
  (only at Γ,q,-q and similar)
\item
  dbg\_do\_smr\_g: metal-only term involving 2 or 3 Fermi-energy shift
  (only at Γ,Γ,Γ)
\item
  dbg\_do\_exc : compute exchange/correlation term
\item
  dbg\_exc\_do\_gga: include gradient correction in the exc term (it is
  always included in the VKS)
\item
  dbg\_add\_core: add core correction to drho (only applies to exc and
  nlcc terms)
\item
  dbg\_do\_nlcc: compute additional non-linear core correction terms
\item
  dbg\_do\_nlcc\_0: nlcc correction to rhod3v
\item
  dbg\_do\_nlcc\_123: nlcc correction to drhod2v
\item
  dbg\_write\_d3\_parts: write each contribution to the D3 matrix
  separately (DEFAULT: .false.)
\item
  dbg\_full\_bands: solve the Sternheimer equation using all the empty
  bands (DEFAULT: .false.)
\end{itemize}

These variables are useful for code debugging, but we do not recommend
using them for any other reason, especially without knowledge of the
code.

\hypertarget{running-d3q}{%
\section{Running D3Q}\label{running-d3q}}

Once you have prepared you d3q.x calculation, running pw.x and ph.x and
your input file is ready, it is time to start d3q.x. This code can take
advantage of MPI parallelism in order to speedup calculation in the same
way as pw.x and ph.x do. On the other hand, it only supports the
plane-waves and k-points pools levels of parallelism. The more advanced
task groups, band groups are not supported. Contrary to phonon it is
also not possible to use image parallelism to run more triplets at the
same time. It is not strictly necessary to run d3q.x with the same
number of CPUs or pools than pw.x, however this is the most commonly
tested case and there is a small chance that some unexpected bug may
appear if you do otherwise. The fildrho files produced by phonon do not
depends on the number of CPUs or pools, it is perfectly safe to reuse
them with different number of CPUs.

\hypertarget{note-on-pools-parallelisation}{%
\subsection{Note on pools
parallelisation}\label{note-on-pools-parallelisation}}

Update: as of QE v6.x, d3q and pw can run with a different number of
pools and pw can be restarted with more pools than k-points.

\hypertarget{output-format}{%
\section{Output format}\label{output-format}}

The d3q.x code will write file containing the D3 matrix for every
triplet it computes and for every triplet that can be obtained from it
using symmetry operations. The name of the file depends on the
coordinates of the q-points in the triplet, but it always start with the
value of the fild3dyn input variable.

The filename is constructed as:

\begin{verbatim}
$fild3dyn_Q1.q11_q12_q13_Q2.q21.q22.q23_Q3.q31.q32.q33
\end{verbatim}

Where the parts in bold are constants and the components qij are is the
j-th coordinate of q point number I, in crystalline fractional
coordinates, expressed as a fraction (i.e.~1/2, 1/3) with the
\enquote{/} replaced with an \enquote{o}. E.g. these are the files
corresponding to the triplets (Γ, Γ, Γ), (Γ, X, -X) and (X, M, -X-M) of
an hexagonal system:

\begin{verbatim}
anh_Q1.0_0_0_Q2.0_0_0_Q3.0_0_0 
anh_Q1.0_0_0_Q2.0_-1o2_0_Q3.0_1o2_0 
anh_Q1.0_-1o2_0_Q2.-1o2_0_0_Q3.1o2_1o2_0
\end{verbatim}

Each file contains the D3 matrix (more strictly, the D3 tensor) in an
XML format that is heavily commented and should be self-documented. At
the beginning of each file, there is a description of the system (cell,
coordinates, masses, etc) then the tensor is saved in Cartesian
coordinates, in units of Ry/bohr3. The tensor is divided in blocks, each
block correspond to three atoms, and is composed of 27=3x3x3 complex
numbers corresponding to the x,y and z directions.

\hypertarget{data-flow}{%
\section{Data flow}\label{data-flow}}

\begin{figure}
\centering
\includegraphics[width=0.65\textwidth,height=\textheight]{images/dataflow.png}
\caption{Data flow from pw to ph to d3q.}
\end{figure}

Note: All the d3\_ codes are documented in manual-thermal2.

\hypertarget{bibliography}{%
\section{Bibliography}\label{bibliography}}

See the \href{thermal2-manual.md\#bibliography}{bibliography section of
the thermal2 manual}.

\end{document}
